% =========================================================
%  sesi/03/sesi.tex
%  Sesi 3 — Aset, Animasi, dan UI Dasar
% =========================================================

\sesi{03}{Aset, Animasi, dan UI Dasar}{4.1}{140 menit}

Ini pertemuan minggu ketiga belas. Level Rak sudah dapat dimainkan, tetapi
seluruh isinya masih berupa balok dan petak berwarna. Sesi ini mengganti
tampilan itu dengan \istilah{sprite} sesungguhnya, menambahkan animasi pada
pemain, dan menampilkan skor serta nyawa pada layar.

Tiga pekerjaan itu mengikuti Bab XIV. Pertama, aset gambar diimpor dan
pengaturan impornya ditetapkan. Kedua, \istilah{Animator} disusun sebagai mesin
keadaan yang memilih animasi diam, berlari, atau melompat. Ketiga,
\istilah{Canvas} dibangun sebagai wadah antarmuka yang menampilkan skor dan
nyawa.

Mekanik yang telah selesai pada sesi 2 tidak diubah. Yang berubah hanya
tampilannya. Karena itu, setiap kali sebuah aset dipasang, jalankan kembali
scene dan pastikan pemain masih dapat berjalan, melompat, dan mengambil
artefak seperti sebelumnya.

% ---------------------------------------------------------
\bagian{Tujuan Sesi}

\begin{tujuanbox}
Setelah menuntaskan sesi ini, Anda dapat menunjukkan:
\begin{enumerate}[leftmargin=1.5em, itemsep=0.2em]
  \item aset \istilah{sprite} yang terimpor dengan \menu{Sprite Mode},
        \menu{Pixels Per Unit}, \menu{Filter Mode}, dan \menu{Compression}
        yang tercatat;
  \item pemain yang gambarnya telah diganti, dengan bentuk batas yang masih
        sesuai dan mekanik yang tidak berubah;
  \item \istilah{Animator} berisi keadaan diam, berlari, dan melompat yang
        berganti menurut gerak pemain;
  \item \istilah{Canvas} yang menampilkan skor dan nyawa, serta berubah ketika
        artefak diambil atau nyawa berkurang;
  \item perbandingan scene sebelum dan sesudah penggantian aset.
\end{enumerate}
\end{tujuanbox}

% ---------------------------------------------------------
\bagian{Prasyarat}

\begin{prasyaratbox}
\begin{itemize}[leftmargin=1.4em, itemsep=0.25em]
  \item Checkpoint sesi 2 sudah berjalan: pemain dapat berjalan, melompat,
        mengambil artefak, dan mencapai pintu.
  \item Anda telah membaca Bab XIV bagian 14.1 sampai 14.5, yaitu bagian
        tentang jalan pengolahan aset, \istilah{Animator}, \istilah{Canvas},
        penggantian \istilah{greybox}, dan lisensi aset.
  \item Anda memiliki aset \istilah{sprite} berlisensi jelas untuk pemain,
        petak, artefak, dan pintu. Catat sumber serta lisensinya sejak awal.
\end{itemize}
\end{prasyaratbox}

% ---------------------------------------------------------
\bagian{Peta Kemampuan}

\petasesi
  {Level yang dapat dimainkan, tetapi seluruhnya masih berupa balok dan petak
   berwarna.}
  {Impor \istilah{sprite}, \istilah{Animator}, dan \istilah{Canvas} berisi
   skor serta nyawa.}
  {Audio, \istilah{trigger} interaksi, dan pengujian prototipe pada sesi 4
   dan 5.}

% ---------------------------------------------------------
\bagian{Konsep yang Dipakai}

\subsection*{Aset diolah ketika diimpor}

Pada Unity, gambar diolah ketika berkasnya masuk ke folder \file{Assets}, bukan
ketika permainan dijalankan. Hasil pengolahan disimpan, sehingga mengubah
pengaturan impor langsung mengubah seluruh benda yang memakai gambar itu.
Karena itu, pengaturan impor adalah tempat pertama yang diperiksa ketika sebuah
aset tampak salah.

Empat pengaturan yang paling berpengaruh: \menu{Sprite Mode} menentukan satu
gambar atau lembar berisi banyak gambar; \menu{Pixels Per Unit} menentukan
berapa piksel yang setara satu satuan; \menu{Filter Mode} menentukan
penghalusan gambar; dan \menu{Compression} menentukan pemadatan gambar.
Ukuran gambar pada scene dihitung dari tinggi berkas dibagi nilai
\menu{Pixels Per Unit}. Seluruh aset sebaiknya memakai nilai yang sama agar
ukurannya dapat dibandingkan.

\subsection*{Animator adalah mesin keadaan}

\istilah{Animator} menyimpan keadaan animasi dan aturan peralihannya. Keadaan
disebut \istilah{animation clip}; peralihannya diatur oleh parameter. Tiga
jenis parameter yang dipakai di sini: angka pecahan untuk kecepatan, nilai
benar atau salah untuk keadaan menyentuh tanah, dan pemicu untuk lompatan.
Pemicu cocok bagi kejadian yang berlangsung sekali, sebab mesin mengembalikan
nilainya sendiri setelah dipakai.

Kode Anda tidak memutar animasi. Kode hanya menetapkan nilai parameter,
sedangkan peralihannya dikerjakan mesin menurut aturan yang Anda gambar. Nama
parameter pada kode harus sama persis dengan nama pada jendela
\istilah{Animator}, termasuk besar-kecil hurufnya. Kesamaan itu tidak
diperiksa, dan kesalahannya tidak menimbulkan pesan apa pun.

\subsection*{Canvas hanya menampilkan}

Elemen antarmuka tidak berada pada ruang permainan dan tidak terpengaruh
kamera. Karena itu, elemen itu ditempatkan pada \istilah{Canvas} tersendiri.
Antarmuka tidak menyimpan keadaan permainan; skor dan nyawa tetap disimpan pada
kode, sedangkan tulisan di layar hanya mengikuti perubahan nilai itu.

% ---------------------------------------------------------
\bagian{Tutorial}

Nama menu dan jendela ditulis seperti yang tampak pada Unity LTS. Contoh pada
sesi ini memakai bingkai berukuran $64 \times 64$ piksel. Apabila aset Anda
berukuran lain, sesuaikan nilai \menu{Pixels Per Unit} menurut hitungan pada
bagian Konsep.

\begin{langkah}
  \item \textbf{Siapkan asetnya.} Unduh aset pemain yang memuat animasi diam,
        berlari, dan melompat, beserta aset petak, artefak, dan pintu. Catat
        nama berkas, sumber, dan lisensi setiap aset pada Lembar Kerja.

  \item \textbf{Buat folder dan impor berkas.} Pada jendela \menu{Project},
        buat folder \code{Sprites} dan \code{Animations} di bawah
        \file{Assets} melalui \menu{Create > Folder}. Salin berkas gambar ke
        \file{Assets/Sprites} melalui \menu{File Explorer}; Unity mengimpornya
        otomatis.

  \item \textbf{Atur impor dan potong lembarnya.} Pilih berkas pemain pada
        jendela \menu{Project}. Pada \menu{Inspector}, setel \menu{Sprite
        Mode} menjadi \menu{Multiple}, \menu{Pixels Per Unit} 64,
        \menu{Filter Mode} \menu{Point}, dan \menu{Compression} \menu{None},
        lalu tekan \menu{Apply}. Tekan \menu{Sprite Editor}, pilih
        \menu{Slice}, setel \menu{Type} menjadi \menu{Grid By Cell Size}
        dengan ukuran $64 \times 64$, lalu tekan \menu{Slice} dan
        \menu{Apply}.

  \item \textbf{Pasang gambar pemain.} Pilih \code{Pemain} pada
        \menu{Hierarchy}. Pada \code{Sprite Renderer}, isi \menu{Sprite}
        dengan salah satu bingkai diam. Periksa \code{Capsule Collider2D}:
        bila gambar barunya lebih besar, sesuaikan \menu{Size} agar tetap
        $(0.7,0.9)$.

  \item \textbf{Tambahkan Animator.} Pilih \code{Pemain}, tekan
        \menu{Add Component}, lalu tambahkan \menu{Animator}. Buat
        \menu{Create > Animator Controller} bernama \code{AnimatorPemain} di
        folder \file{Assets/Animations}, seret ke kolom \menu{Controller},
        lalu klik dua kali untuk membuka jendela
        \menu{Window > Animation > Animator}.

  \item \textbf{Buat tiga keadaan.} Pilih seluruh bingkai diam dengan
        \menu{Ctrl+klik} pada jendela \menu{Project}, lalu seret ke jendela
        \menu{Animator} dan beri nama \code{Idle}. Klik kanan keadaannya,
        pilih \menu{Set as Layer Default State}. Ulangi untuk bingkai
        berlari dengan nama \code{Run} dan bingkai melompat dengan nama
        \code{Jump}.

  \item \textbf{Buat parameter.} Pada tab \menu{Parameters}, tekan \menu{+}
        tiga kali: \menu{Float} bernama \code{speed}, \menu{Bool} bernama
        \code{grounded}, dan \menu{Trigger} bernama \code{lompat}. Huruf
        kecil pada nama itu harus sama dengan yang dipakai skrip.

  \item \textbf{Transisi diam dan berlari.} Klik kanan \code{Idle}, pilih
        \menu{Make Transition}, lalu klik \code{Run}. Pada
        \menu{Inspector} transisi, matikan \menu{Has Exit Time}, setel
        \menu{Transition Duration} 0, dan tambahkan syarat
        \code{speed} \menu{Greater} 0.1. Buat transisi balik
        \code{Run} ke \code{Idle} dengan syarat \code{speed} \menu{Less} 0.1.

  \item \textbf{Transisi lompat.} Klik kanan \code{Any State}, pilih
        \menu{Make Transition}, lalu klik \code{Jump}. Matikan
        \menu{Has Exit Time} dan tambahkan syarat \code{lompat}. Buat
        transisi \code{Jump} ke \code{Idle} dengan syarat
        \code{grounded} bernilai benar.

  \item \textbf{Pasang skrip animasi.} Seret \code{AnimasiPemain} ke
        \code{Pemain}, lalu isi \menu{Body} dengan komponen
        \code{Rigidbody2D} pemain. Skrip ini membalik arah hadap gambar
        menurut arah geraknya.

  \item \textbf{Perbarui KontrolPemain.} Ganti isi
        \file{Scripts/KontrolPemain.cs} dengan versi pada bagian Kode Sesi
        Ini. Pada \menu{Inspector}, isi kolom \menu{Animator} dengan komponen
        \menu{Animator} pemain, lalu jalankan scene.
\end{langkah}

\begin{ujibox}
Setelah Langkah 11, periksa empat hal berikut.
\begin{itemize}[leftmargin=1.4em, itemsep=0.2em]
  \item Pemain menampilkan gambar diam ketika berhenti.
  \item Animasi berlari muncul ketika pemain bergerak, dan kembali diam ketika
        berhenti.
  \item Animasi melompat muncul sekali setiap lompatan, bukan berulang selama
        tombol ditahan.
  \item Gambar pemain menghadap ke kiri ketika bergerak ke kiri, dan ke kanan
        ketika bergerak ke kanan.
\end{itemize}
\end{ujibox}

\begin{langkah}
  \item \textbf{Buat Canvas.} Pilih menu
        \menu{GameObject > UI > Canvas}. Unity menambahkan \code{Canvas} dan
        \code{EventSystem} pada \menu{Hierarchy}. Pada \menu{Canvas Scaler},
        setel \menu{UI Scale Mode} menjadi \menu{Scale With Screen Size},
        \menu{Reference Resolution} $1920 \times 1080$, dan \menu{Match} 0.5.

  \item \textbf{Buat teks skor dan nyawa.} Klik kanan \code{Canvas}, pilih
        \menu{UI > Legacy > Text}, dan beri nama \code{TextSkor}. Setel
        \menu{Anchor} kiri atas, \menu{Pos} $(30,-30)$, \menu{Size}
        $400 \times 60$, \menu{Font Size} 36, \menu{Color} putih, dan
        \menu{Text} \code{Skor: 0}. Gandakan teks itu, beri nama
        \code{TextNyawa}, setel \menu{Anchor} kanan atas, \menu{Pos}
        $(-30,-30)$, \menu{Width} 300, \menu{Alignment} kanan, dan
        \menu{Text} \code{Nyawa: 3}.

  \item \textbf{Buat ikon nyawa.} Klik kanan \code{Canvas}, pilih
        \menu{UI > Image}, dan beri nama \code{IkonNyawa1}. Setel
        \menu{Anchor} kanan atas, \menu{Size} $32 \times 32$, \menu{Pos}
        $(-30,-80)$, dan \menu{Color} \code{#D94A4A}. Gandakan dua kali
        menjadi \code{IkonNyawa2} dan \code{IkonNyawa3}, lalu geser
        \menu{Pos X} masing-masing menjadi $-70$ dan $-110$.

  \item \textbf{Pasang skrip HUD.} Seret \code{TampilanHud} ke
        \code{Canvas}. Pada \menu{Inspector}, isi \menu{Score Text} dengan
        \code{TextSkor}, \menu{Lives Text} dengan \code{TextNyawa}, dan
        \menu{Heart Icons} dengan ukuran 3 berisi ketiga ikon nyawa.

  \item \textbf{Perbarui Artefak.} Ganti isi \file{Scripts/Artefak.cs}
        dengan versi pada bagian Kode Sesi Ini. Pada setiap artefak, isi
        kolom \menu{Hud} dengan objek \code{Canvas}.

  \item \textbf{Ganti aset lingkungan.} Ganti gambar petak lantai dan rak
        melalui \menu{Sprite} pada aset \code{TileLantai} dan
        \code{TileRak}. Ganti pula gambar artefak dan pintu. Periksa kembali
        \code{Collider} setiap benda, sebab bentuk batas tidak berubah
        otomatis bersama gambarnya.

  \item \textbf{Bandingkan hasilnya.} Simpan tangkapan layar scene sebelum
        dan sesudah penggantian aset. Jalankan scene, ambil ketiga artefak,
        lalu uji nyawa melalui klik kanan komponen \code{TampilanHud} dan
        pilih \menu{Uji kurangi satu nyawa}. Simpan scene dengan
        \menu{Ctrl+S}.
\end{langkah}

\begin{ujibox}
Setelah Langkah 18, periksa lima hal berikut.
\begin{itemize}[leftmargin=1.4em, itemsep=0.2em]
  \item Seluruh benda memakai gambar aset, bukan lagi kotak berwarna.
  \item \menu{TextSkor} berubah dari \code{Skor: 0} menjadi
        \code{Skor: 3} setelah ketiga artefak diambil.
  \item \menu{TextNyawa} dan ikon nyawa berkurang ketika menu konteks
        \menu{Uji kurangi satu nyawa} dipilih.
  \item Pemain masih dapat berjalan, melompat, dan mendarat seperti pada
        sesi 2.
  \item Tidak ada pesan merah pada \menu{Console}.
\end{itemize}
\end{ujibox}

% ---------------------------------------------------------
\bagian{Kode Sesi Ini}

\code{KontrolPemain} bertambah bagian animasi: ia menetapkan kecepatan dan
keadaan tanah, serta memicu parameter lompatan. \code{AnimasiPemain}
membalik arah hadap gambar. \code{Artefak} memanggil fungsi skor pada HUD.
\code{TampilanHud} menyimpan skor dan nyawa, lalu memperbarui tampilannya.
\code{Pintu} tidak berubah dari sesi 2; salinannya tetap disertakan pada
checkpoint.

\tampilkode[language=CSharp]{code/sesi-03/Scripts/KontrolPemain.cs}
  {KontrolPemain.cs: kontrol pemain yang juga menetapkan parameter Animator}

\tampilkode[language=CSharp]{code/sesi-03/Scripts/AnimasiPemain.cs}
  {AnimasiPemain.cs: membalik arah hadap gambar menurut arah gerak}

\tampilkode[language=CSharp]{code/sesi-03/Scripts/Artefak.cs}
  {Artefak.cs: trigger yang menambah skor melalui TampilanHud}

\tampilkode[language=CSharp]{code/sesi-03/Scripts/TampilanHud.cs}
  {TampilanHud.cs: menampilkan skor dan nyawa pada Canvas}

\subsection*{Tabel objek scene}

Tabel berikut adalah daftar objek yang harus ada pada scene Anda. Nilai yang
tidak disebut boleh dibiarkan bawaan.

{\footnotesize
\renewcommand{\arraystretch}{1.2}
\begin{tabularx}{\linewidth}{@{}L{2.3cm} L{3.5cm} Y@{}}
\toprule
\textbf{GameObject} & \textbf{Komponen} & \textbf{Nilai penting} \\
\midrule
Pemain & Sprite Renderer, Rigidbody2D, Capsule Collider2D, Animator,
  KontrolPemain, AnimasiPemain &
  Animator Controller \code{AnimatorPemain}; parameter \code{speed},
  \code{grounded}, \code{lompat}; Body \code{Rigidbody2D}; Ambang Gerak 0.1 \\
Canvas & Canvas, Canvas Scaler, Graphic Raycaster, TampilanHud &
  Screen Space Overlay; Scale With Screen Size; Reference
  $1920 \times 1080$; Match 0.5 \\
TextSkor & Rect Transform, Text &
  Anchor kiri atas; Pos $(30,-30)$; $400 \times 60$; Font Size 36; putih;
  \code{Skor: 0} \\
TextNyawa & Rect Transform, Text &
  Anchor kanan atas; Pos $(-30,-30)$; $300 \times 60$; Font Size 36; putih;
  \code{Nyawa: 3} \\
IkonNyawa1--3 & Rect Transform, Image &
  Anchor kanan atas; $32 \times 32$; Pos X $-30$, $-70$, $-110$; Pos Y $-80$;
  \code{#D94A4A} \\
Artefak1--3 & Sprite Renderer, Circle Collider2D, Artefak &
  Is Trigger; Radius 0.25; Nilai 1; Hud berisi \code{Canvas} \\
Pintu & Sprite Renderer, Box Collider2D, Pintu &
  $(8.3,-2.6,0)$; $(1,1.8)$; Is Trigger; Artefak Diperlukan 3 \\
\bottomrule
\end{tabularx}
}

\subsection*{Tabel pengaturan impor}

{\footnotesize
\renewcommand{\arraystretch}{1.2}
\begin{tabularx}{\linewidth}{@{}L{3.4cm} L{2.3cm} C{1.6cm} L{1.9cm} Y@{}}
\toprule
\textbf{Aset} & \textbf{Sprite Mode} & \textbf{PPU} & \textbf{Filter} &
\textbf{Compression} \\
\midrule
Pemain (lembar animasi) & Multiple & 64 & Point & None \\
Petak lantai dan rak & Single & 64 & Point & None \\
Artefak dan pintu & Single & 64 & Point & None \\
\bottomrule
\end{tabularx}
}

\checkpoint{03}{\raggedright Folder checkpoint memuat \file{Scripts/KontrolPemain.cs},
\file{Scripts/Artefak.cs}, \file{Scripts/Pintu.cs},
\file{Scripts/AnimasiPemain.cs}, \file{Scripts/TampilanHud.cs}, dan
\file{scene.md}. Scene dan berkas aset tidak disertakan sebagai berkas biner;
susunlah mengikuti tabel di atas.}

% ---------------------------------------------------------
\bagian{Kesalahan Umum}

\begin{galatbox}
\begin{itemize}[leftmargin=1.4em, itemsep=0.35em]
  \item \textbf{Animasi tidak berganti atau transisi berulang.}
        Gejala: gambar tetap diam, melompat terus-menerus, atau kembali ke
        diam di tengah lompatan, tanpa pesan kesalahan. Periksa nama parameter
        pada jendela \istilah{Animator} terhadap nama pada skrip, nilai
        parameter ketika scene berjalan, dan \menu{Has Exit Time} pada
        transisi.
  \item \textbf{Ukuran, arah, atau bentuk batas tidak sesuai.}
        Gejala: pemain raksasa, menghadap ke belakang, atau bertabrakan
        seperti balok lama. Periksa \menu{Pixels Per Unit}, skala $x$ pada
        \code{AnimasiPemain}, serta \menu{Size} dan \menu{Offset} pada
        \code{Capsule Collider2D}. Mengganti \menu{Sprite} tidak mengubah
        collider.
  \item \textbf{Skor atau nyawa tidak berubah.}
        Gejala: tulisan tetap \code{Skor: 0} dan Console menampilkan
        peringatan. Periksa rujukan \menu{Score Text}, \menu{Lives Text}, dan
        \menu{Hud} pada artefak. Rujukan yang kosong membuat skrip HUD
        mematikan dirinya sendiri.
\end{itemize}
\end{galatbox}

% ---------------------------------------------------------
\bagian{Lembar Kerja}

Isi tabel berikut dan simpan pada catatan praktikum Anda.

{\small
\renewcommand{\arraystretch}{1.4}
\begin{tabularx}{\linewidth}{@{}C{0.8cm} L{6.2cm} Y@{}}
\toprule
\textbf{No} & \textbf{Hasil yang dicatat} & \textbf{Keterangan Anda} \\
\midrule
1 & Nama, sumber, dan lisensi aset pemain & \\[0.45cm]
2 & Nama, sumber, dan lisensi aset petak, artefak, dan pintu & \\[0.45cm]
3 & Nilai Pixels Per Unit dan alasan pemilihannya & \\[0.45cm]
4 & Jumlah bingkai pada setiap animation clip & \\[0.45cm]
5 & Nilai ambang kecepatan pada transisi diam--berlari & \\[0.45cm]
6 & Perubahan mekanik yang timbul setelah aset diganti & \\[0.7cm]
7 & Masalah animasi atau UI yang ditemui dan perbaikannya & \\[0.7cm]
\bottomrule
\end{tabularx}
}

% ---------------------------------------------------------
\bagian{Checklist Sesi}

\begin{checklistbox}
\begin{checklist}
  \item Saya dapat menunjukkan aset \istilah{sprite} yang terimpor dengan
        pengaturan seragam, dan menghitung tinggi gambarnya pada scene.
  \item Saya dapat menunjukkan \istilah{Animator} dengan keadaan diam,
        berlari, dan melompat beserta transisinya.
  \item Saya dapat menjelaskan perbedaan parameter angka, nilai benar-salah,
        dan pemicu, serta memilih jenis yang tepat untuk lompatan.
  \item Saya dapat menunjukkan \istilah{Canvas} yang skor dan nyawanya
        berubah mengikuti keadaan permainan.
  \item Saya dapat membandingkan scene sebelum dan sesudah penggantian aset.
\end{checklist}
\end{checklistbox}

% ---------------------------------------------------------
\bagian{Pengait ke Sesi Berikutnya}

Tampilan level kini lengkap, tetapi masih bisu. Sesi 4 menambahkan bunyi latar
dan efek bunyi, serta menghubungkan lompat, pengambilan artefak, dan pintu
dengan \istilah{trigger}.
